% !TEX root = ../main.tex
\section{The Hypergraph Kernel}\label{sec:kernel}

Barkeshli, Douglas, and Freedman~\cite{BDF2026} model mathematical reasoning as a directed hypergraph whose hyperedges connect $p$ inputs to $q$ outputs.  We instantiate their framework on CIC by defining two primitives: a \emph{nerve} (a single content-addressed hyperedge) and an \emph{AstroNet} (a finite set of nerves).  The kernel operates exclusively on these two structures.

\begin{definition}[AstroNerve]\leanlink{https://github.com/MathNetwork/cic-kernal/blob/main/HypergraphKernel/HypergraphKernel/AstroNerve.lean\#L21}
An \textbf{AstroNerve} is a triple $(\mathrm{id},\, \mathrm{ref},\, \mathrm{record})$ where $\mathrm{id}$ is a hash, $\mathrm{ref} = (r_1, \ldots, r_k)$ is an ordered list of reference hashes, and $\mathrm{record}$ is a serialized content string.  The \textbf{width} is $\omega(e) = |\mathrm{ref}(e)| - 1$.  Nerves with $\omega = 0$ are vertices; nerves with $\omega \geq 1$ are hyperedges.
\end{definition}

\begin{definition}[AstroNet]\label{def:ca-hypergraph}\label{def:well-formed}\leanlink{https://github.com/MathNetwork/cic-kernal/blob/main/Theory/Theory/Core/Basic.lean\#L20}
An \textbf{AstroNet} is a finite set $A$ of nerves together with a hash function $H$, satisfying five conditions:
\begin{enumerate}
  \setcounter{enumi}{-1}
  \item \textbf{Content-addressing.} $\mathrm{id}(e) = H(\mathrm{record}(e))$ for every $e \in A$.
  \item \textbf{Self-reference.} Every atom ($\omega(e) = 0$) satisfies $\mathrm{ref}(e) = (\mathrm{id}(e))$.
  \item \textbf{Injectivity.} $e \mapsto \mathrm{id}(e)$ is injective on $A$.
  \item \textbf{Closure.} For every $r_i \in \mathrm{ref}(e)$, there exists $e' \in A$ with $\mathrm{id}(e') = r_i$.
  \item \textbf{No self-reference.} For $\omega(e) > 0$, $\mathrm{id}(e) \notin \mathrm{ref}(e)$.
\end{enumerate}
The $\mathrm{ref}$ list may contain repeated hashes (e.g., $\mathrm{ref} = [h_{\mathrm{Nat}}, h_{\mathrm{Nat}}]$ for $\texttt{forallE}(\texttt{Nat}, \texttt{Nat})$).
\end{definition}

These conditions support a depth filtration formalized in the Theory module (534 lines, using Mathlib's \texttt{Finset} library).

\begin{definition}[Depth filtration]\leanlink{https://github.com/MathNetwork/cic-kernal/blob/main/Theory/Theory/Core/Basic.lean\#L76}
Given an AstroNet $A$, define the filtration $A^{(0)} \subseteq A^{(1)} \subseteq \cdots$ by
\[
  A^{(0)} = \{ e \in A \mid \omega(e) = 0 \}, \qquad
  A^{(n+1)} = \{ e \in A \mid \forall\, r \in \mathrm{ref}(e),\; r \in \mathrm{ids}(A^{(n)}) \}.
\]
The \textbf{depth} of $e$ is $\delta(e) = \min\{n \mid e \in A^{(n)}\}$.  Nerves on reference cycles have no finite depth.
\end{definition}

This filtration is monotone and stabilizes on any finite AstroNet.  Staged nerves are acyclic; unstaged nerves reach cycles.  Proofs appear in Appendix~\ref{app:proofs}.

In the construction of \texttt{Nat.add\_comm} and its 14 prerequisites, the hypergraph grows from empty to 612 nerves: 15 declaration nerves and 597 expression nerves.

Figure~\ref{fig:succ-compare} illustrates the difference on a simple example: the type of \texttt{Nat.succ}, which is $\mathbb{N} \to \mathbb{N}$ (\texttt{forallE(Nat, Nat)}).

\begin{figure}[H]
\centering
\begin{subfigure}[t]{0.44\textwidth}
\centering
\begin{tikzpicture}[
    nd/.style={draw, rounded corners=3pt, fill=gray!10, inner sep=3pt, font=\small\ttfamily},
    dup/.style={draw, rounded corners=3pt, fill=red!10, thick, inner sep=3pt, font=\small\ttfamily},
    edge/.style={-Latex, thin},
    lbl/.style={font=\scriptsize, gray}
]
  \node[nd] (root) at (0, 0) {forallE};
  \node[dup] (nat1) at (-1.5, -1.5) {Nat};
  \node[dup] (nat2) at (1.5, -1.5) {Nat};
  \draw[edge] (root) -- (nat1);
  \draw[edge] (root) -- (nat2);
  \node[lbl] at (-1.5, -2.2) {copy 1};
  \node[lbl] at (1.5, -2.2) {copy 2};
\end{tikzpicture}
\caption{Expression tree: 3 nodes, two copies of \texttt{Nat}.}
\label{fig:succ-tree}
\end{subfigure}\hfill
\begin{subfigure}[t]{0.44\textwidth}
\centering
\begin{tikzpicture}[
    nd/.style={draw, rounded corners=3pt, fill=gray!10, inner sep=3pt, font=\small\ttfamily},
    shared/.style={draw, rounded corners=3pt, fill=green!15, thick, inner sep=3pt, font=\small\ttfamily},
    edge/.style={-Latex, thin},
    lbl/.style={font=\scriptsize, gray}
]
  \node[nd] (root) at (0, 0) {forallE};
  \node[shared] (nat) at (0, -1.5) {Nat};
  \draw[edge] (root) to[out=-150, in=90] (-0.3, -0.9) -- (nat);
  \draw[edge] (root) to[out=-30, in=90] (0.3, -0.9) -- (nat);
  \node[lbl] at (0, -2.2) {one nerve, two refs};
  \node[lbl] at (0, -2.8) {ref $= [h_{\mathrm{Nat}}, h_{\mathrm{Nat}}]$};
\end{tikzpicture}
\caption{Hypergraph: 2 nerves, \texttt{Nat} shared.}
\label{fig:succ-hyper}
\end{subfigure}
\caption{The type of \texttt{Nat.succ} ($\mathbb{N} \to \mathbb{N}$) in the two kernels.  Left: two independent \texttt{Nat} nodes.  Right: one \texttt{Nat} nerve referenced twice.}
\label{fig:succ-compare}
\end{figure}
